\documentclass[11pt]{article}

\usepackage[margin=1.2in]{geometry}
\usepackage{amsmath,amssymb}
\usepackage{microtype}
\usepackage{parskip}
\usepackage{hyperref}
\usepackage{xcolor}

\title{The FEN Schema Security Architecture}
\author{M. A. Thomas}
\date{}

\begin{document}

\maketitle

\section{Core Principle}

The FEN schema operates at the semantic level. Its objects---Facts,
Episodes, Narratives---are abstracted from any particular storage
technology. The security model must follow suit. The primitive is not
``encrypt this PDF'' or ``encrypt this database row.'' The primitive is
``protect this hemoglobin Fact,'' ``protect this oncology Episode,''
``protect this patient-authored Narrative.''

This matters because FEN objects do not map one-to-one onto storage
objects. A single Episode may draw from many stored records. A single
Fact may have provenance in one location, payload data in another, and
policy metadata in a third. Tying security to the storage model produces
a brittle system whose protection boundaries diverge from its semantic
boundaries.

\paragraph{Three levels of access control.}

\begin{enumerate}
\item \textbf{Storage access.} The ability to fetch encrypted blobs and
  ciphertext from the infrastructure. Unavoidable for any process that
  touches the storage layer.

\item \textbf{Cryptographic access.} The ability to obtain decryption
  keys. Governed by policy; scoped to specific principals and purposes.

\item \textbf{Semantic access.} The ability to reconstruct a plaintext
  Fact, Episode, or Narrative as a meaningful object. Requires both
  cryptographic access and satisfaction of the governing consent policy.
\end{enumerate}

The security layer governs semantic access directly and shapes
cryptographic access. Its relation to storage access depends on
infrastructure choices.

\paragraph{Derivations.}
Many useful views are assembled on demand from underlying Facts and
Episodes. Once we allow recomposition, encryption policy must govern
derived objects too. A generated Narrative, timeline, or export inherits
constraints from its constituents. The composition rules are specified in
Section~\ref{sec:consent-composition}.

\section{Threat Model}
\label{sec:threat-model}

The architecture must address the following adversary scenarios. They are
ordered roughly by likelihood in a production health-data platform.

\begin{enumerate}

\item \textbf{Curious employee.} A Phoros engineer or operator with
  access to operational infrastructure queries \texttt{FactOperational}
  records. They should not be able to decrypt payloads without triggering
  an audit event. This is the most common real-world breach vector in
  health technology.

\item \textbf{Compromised service account.} An automated process with
  legitimate access to some encrypted payloads has its credentials
  stolen or misused. Blast radius depends on key scoping---per-person,
  per-Fact-type, or platform-wide.

\item \textbf{Malicious data borrower.} A party with authorized access
  to a cohort's data attempts to re-identify individuals, retain data
  beyond the authorized period, or combine disclosed data with external
  datasets. This is a contractual and cryptographic problem
  simultaneously.

\item \textbf{Infrastructure compromise.} An attacker gains access to
  the storage layer---databases, object stores, backups. They see
  ciphertext and \texttt{FactOperational} metadata.
  Section~\ref{sec:factoperational-leakage} addresses what this reveals.

\item \textbf{Colluding insiders.} Two or more Phoros employees, or one
  employee and one external party, combine their access to exceed what
  either should have alone. Multi-party key schemes mitigate this.

\item \textbf{Legal compulsion.} A subpoena or court order demands
  patient data. The architecture should make it possible to comply
  narrowly---producing specific Facts for a specific person---without
  exposing the broader dataset. Granular encryption supports this.

\item \textbf{Phoros itself as adversary.} The patient's data should be
  protected even if Phoros as an organization acts against the patient's
  interests. This is aspirational and probably not fully achievable in
  v1, but stating it clarifies where the architecture falls short and
  where future hardening should focus.

\end{enumerate}

Key management, including key derivation hierarchy, rotation, and
compromise recovery, is deferred to a separate document. The present
document assumes that keys are policy-governed and scoped, without
specifying the mechanism.

\section{Eager Materialization}

There are two senses in which FEN objects get materialized.
\emph{Semantic materialization} means constructing FEN objects from raw
inputs (e.g., data ingested from Epic). \emph{Plaintext materialization}
means decrypting data so a principal can read it. Since everything is
encrypted at rest, plaintext materialization happens every time someone
accesses data. The eager-vs-lazy question concerns semantic
materialization.

We adopt eager semantic materialization as the default. Incoming source
data are translated into canonical FEN objects at ingestion time. The
system stores encrypted Facts, Episodes, and other semantic units as the
authoritative protected objects, each with a stable identity and
associated policy metadata.

The advantage: the object being encrypted and the object being governed
are the same object. Access control, auditing, provenance binding, and
consent enforcement all attach to entities that already carry meaning in
the system.

The system permits deferred construction only for higher-order derived
objects---Episodes or Narratives generated from existing Facts. In these
cases, the derived objects inherit protection constraints from their
constituents rather than introducing ungoverned semantic material.

\section{Facts as the Cryptographic Unit}

The central granularity question: what is the smallest unit that can
carry independent provenance, consent constraints, and disclosure
meaning?

The answer is the Fact. A Fact represents an assertion, measurement, or
observation arising from a sufficiently unified act of generation. A lab
result, a vital-sign measurement, a diagnosis assertion, a medication
administration---each corresponds to a single generating event.

Two criteria determine Fact boundaries:

\begin{enumerate}
\item \textbf{Unified generation.} The fields of the Fact are jointly
  produced by a single clinical measurement, observation, or inference
  process.

\item \textbf{Consent separability.} It is sensible to treat the Fact
  as a single object of disclosure or restriction under a patient's
  consent policy. If a component would require independent disclosure
  rules, it should be a separate Fact.
\end{enumerate}

This produces a natural alignment:

\begin{quote}
Facts are the cryptographic unit. Episodes are the clinical unit.
Narratives are the interpretive unit.
\end{quote}

A Fact is encrypted as one payload. Its internal fields stay together
because their meaning arises from the unified event that produced them.
Encrypting fields separately would fragment the semantic object while
multiplying key management complexity.

Episodes and Narratives are protected indirectly through their
constituent Facts. An Episode inherits protection constraints from its
Facts. A Narrative is governed through the combined constraints of its
underlying Facts and Episodes. When the interpretive structure itself
carries sensitive meaning beyond the underlying data, the higher-level
object may receive additional cryptographic wrapping.

\paragraph{Qualifications.}

Certain metadata must remain visible outside the encrypted payload to
support indexing, routing, and policy evaluation. This is the
\texttt{FactOperational} layer described in Section~\ref{sec:factoperational}.

When a candidate Fact contains components with materially different
consent implications, the correct response is usually to model them as
separate Facts, not to encrypt fields separately within one Fact.

Selective disclosure within a Fact (nested protection, partial payload
encryption) should remain exceptional.

\section{Plaintext Materialization}

Facts do not exist in readable form in the storage layer. The readable
object arises when the system reconstructs it through a controlled
procedure: evaluate policy, retrieve encrypted fragments, verify
integrity, decrypt. Plaintext Facts exist only transiently in memory
during authorized access.

Storing or caching plaintext objects in a database would break the
alignment between the semantic model and the security model. We do not
do this.

\paragraph{The procedure.}
Every Fact has a persistent \texttt{FactId}. The corresponding
\texttt{FactOperational} record (Section~\ref{sec:factoperational})
contains what is needed to reconstruct the object: the payload locator,
integrity hash, and policy reference. It is a registry entry, not a
container for plaintext.

When a principal requests a Fact:

\begin{enumerate}
\item The system retrieves the \texttt{FactOperational} record for the
  requested \texttt{FactId}.
\item The policy engine evaluates the request against the governing
  consent conditions, the principal's identity, and the declared purpose
  of use.
\item If policy evaluation succeeds, the system retrieves the encrypted
  payload fragments referenced by the \texttt{PayloadLocator}.
\item Integrity checks confirm the fragments match the stored
  \texttt{integrity\_hash}.
  \textcolor{blue}{What is being hashed (encrypted payload bytes, the semantic Fact fields, both)?  What happens to hash after key rotation if payload is re-encrypted?}
\item Cryptographic keys are released under the governing policy.
\item Decryption produces the payload, which the system assembles into
  the typed Fact defined by the FEN schema.
\end{enumerate}

Episodes and Narratives follow the same principle. Their
\texttt{FactOperational} records (or lightweight operational
equivalents) reference multiple constituent Facts. Materialization
involves recursive reconstruction of constituents, each subject to its
own policy. The resulting object inherits the combined constraints
(Section~\ref{sec:consent-composition}).

\paragraph{Two indexes.}
The architecture separates two forms of indexing. An \emph{association
index} maps a \texttt{PersonId} to associated \texttt{FactId},
\texttt{EpisodeId}, and \texttt{NarrativeId} values. It supports
navigation but contains no clinical payloads. The
\texttt{FactOperational} records themselves serve as the second index:
given a \texttt{FactId}, they provide the reconstruction path to the
encrypted payload.

\section{FactOperational}
\label{sec:factoperational}

The platform needs to index, search, discover cohorts, and evaluate
policy without decrypting payloads. \texttt{FactOperational} is the
operational representation that supports this.

It contains: the stable \texttt{FactId}, the associated
\texttt{PersonId}, Fact type, temporal boundaries, coarse classification
tags, a reference to the governing policy, a locator for encrypted
payload fragments, and integrity verification material. It may also
contain derived indexing features or search tokens for cohort discovery.

\texttt{FactOperational} is not a subset of the semantic Fact. Some
elements (identifiers, temporal boundaries) appear in both. Others
(payload locators, indexing metadata) exist only in the operational
layer. The full semantic representation contains the decrypted payload
that the operational structure intentionally omits.

Different payload types require different operational variants. A
\texttt{VitalsOperational} variant exposes flags like
\texttt{has\_blood\_pressure} and \texttt{abnormal\_flag}. A
\texttt{LabOperational} variant exposes \texttt{lab\_test\_code} and a
bucketed value range. The operational layer mirrors the
\texttt{FactPayload} taxonomy while excluding decrypted clinical values.

\subsection{What FactOperational reveals}
\label{sec:factoperational-leakage}

\texttt{FactOperational} is a selective mitigation of PHI exposure, not
a complete elimination. A Phoros engineer debugging an indexing pipeline
sees \texttt{abnormal\_flag: true} on a lab Fact, not ``creatinine 4.2
mg/dL.'' That is a meaningful reduction in exposure surface for routine
operations. It is not zero-knowledge.

The reduction is uneven across Fact types. \texttt{VitalsOperational}
with \texttt{abnormal\_flag} is relatively low-risk.
\texttt{LabOperational} with \texttt{lab\_test\_code = HIV\_VIRAL\_LOAD}
is high-risk even without a value, because the existence of the test is
itself sensitive. The existence of an oncology-typed Episode is
sensitive. Operational variant design must make sensitivity-aware
decisions about which fields to expose

\textcolor{blue}{Phoros engineers would have access to all FactOperational fields, both low and high risk? Is there a role-based access layer within FactOperational, or is the design assumption that all operational metadata is available to any authenticated Phoros process?}

Access to \texttt{FactOperational} is governed and audited. The fact
that payloads are encrypted does not make operational metadata a
free-for-all. Access controls on the operational layer are less
restrictive than payload decryption, but they exist.

External parties such as data borrowers do not receive access to
\texttt{FactOperational} fields. Borrowers interact only with aggregate
cohort counts until individuals grant explicit permission.

\textcolor{blue}{Idea: add sesitivity tier to FactOperational}

\begin{verbatim}
pub enum OperationalSensitivityTier {
    Routine,       // abnormal flags, visit counts, generic procedure types
    Elevated,      // specific lab test codes, medication classes
    Restricted,    // HIV, oncology, mental health, substance use codes
}
\end{verbatim}


\subsection{Layered materialization}

All PHI is stored encrypted at rest. The security layer permits two
levels of materialization:

\begin{enumerate}
\item \textbf{Operational materialization.} The \texttt{FactOperational}
  record is accessible to authorized Phoros platform processes.
  Operational PHI required for indexing, policy evaluation, and cohort
  discovery is available. The sensitive payload within the
  \texttt{FactPayload} variant remains encrypted.

\item \textbf{Full plaintext materialization.} The complete Fact,
  including payload contents, is reconstructed. Permitted only for the
  individual account holder or for an authorized recipient acting under
  permissions granted by that individual.
\end{enumerate}

\subsection{Interaction with the query optimizer}

The FEN schema's query optimizer
operates over \texttt{FactOperational}, not over decrypted Facts. This
constrains what queries are expressible without decryption. Subject
filtering, time-range filtering, Fact-type slicing, and code-based
narrowing all operate on operational fields. Predicate evaluation that
requires payload content (e.g., ``creatinine above 2.0'') requires
plaintext materialization of candidate Facts and can occur only under
policy authorization.

\subsection{Rust specification}

\begin{verbatim}
use uuid::Uuid;
use chrono::{DateTime, Utc};

pub type FactId = Uuid;
pub type PersonId = Uuid;
pub type PolicyId = Uuid;

#[derive(Debug, Clone)]
pub struct PayloadLocator {
    pub storage_system: String,
    pub object_key: String,
    pub fragment_ids: Vec<String>,
}

#[derive(Debug, Clone)]
pub struct FactOperationalHeader {
    pub fact_id: FactId,
    pub person_id: PersonId,
    pub interval_start: DateTime<Utc>,
    pub interval_end: DateTime<Utc>,
    pub policy_ref: PolicyId,
    pub integrity_hash: String,
    pub payload_locator: PayloadLocator,
    OperationalSensitivityTier // new
}

#[derive(Debug, Clone)]
pub struct VitalsOperational {
    pub header: FactOperationalHeader,
    pub has_blood_pressure: bool,
    pub has_heart_rate: bool,
    pub abnormal_flag: bool,
    pub classification_tags: Vec<String>,
}

#[derive(Debug, Clone)]
pub struct LabOperational {
    pub header: FactOperationalHeader,
    pub lab_test_code: String,
    pub value_bucket: Option<String>,
    pub abnormal_flag: bool,
}

#[derive(Debug, Clone)]
pub enum FactOperational {
    Vitals(VitalsOperational),
    Lab(LabOperational),
    // Medication, Imaging, Procedure, etc.
}
\end{verbatim}

\subsection{Operational representations of Episodes and Narratives}

The case for operational twins is strongest at the Fact level, because
Facts are the primary units of ingestion, indexing, and cryptographic
protection.

\textcolor{blue}{operational twins?}

Episodes can usually be handled through relationships among constituent
Facts plus lightweight metadata (identifier, time interval, episode
type). A full operational twin analogous to \texttt{FactOperational} may
be warranted if episode-level constructs become direct targets of cohort
discovery.

Narratives are the weakest candidate for operational twins. They are
interpretive objects whose value resides almost entirely in their
content.

\section{Consent Composition}
\label{sec:consent-composition}

A derived object---an Episode, a Narrative, a timeline view, an
export---draws on multiple constituent Facts. Each constituent may be
governed by a different consent policy. The system must compute a
resulting policy for the derived object.

\paragraph{Default: intersection.}
The derived object can be disclosed only under conditions that satisfy
every constituent's policy. If a Narrative references Facts $F_1$, $F_2$,
$F_3$ governed by policies $P_1$, $P_2$, $P_3$, then the Narrative is
governed by $P_1 \cap P_2 \cap P_3$.

Example: $P_1$ allows research use by academic institutions. $P_2$
allows research use by any party. $P_3$ allows clinical use only. Under
intersection, the Narrative cannot be disclosed for research, because
$P_3$ forbids it.

\textcolor{blue}{What about partial Episode / Narrative. Could contents be restricted to only F2 or F1+F2? Patient couls allow partial reconstruction?}

Union is never the right default. It would mean the least restrictive
constituent governs the whole, violating the patient's intent on
restricted Facts.

\paragraph{Implications.}

\begin{enumerate}

\item \textbf{One restrictive Fact can lock a derived object.} A single
  highly-restricted Fact makes the entire Episode or Narrative
  inaccessible for otherwise-permitted purposes. The system should
  therefore support constructing derived objects that exclude specific
  Facts when the requesting purpose does not satisfy those Facts'
  policies. A borrower might receive a Narrative built from $F_1$ and
  $F_2$ but not $F_3$.

\item \textbf{Policy granularity affects composition cost.} If every
  Fact has a bespoke policy, intersection is expensive to compute and
  hard to reason about. The recommendation is a small vocabulary of
  policy templates (e.g., ``research-ok,'' ``clinical-only,''
  ``fully-restricted'') with per-Fact overrides in exceptional cases.

\item \textbf{Composition must precede materialization.} The system must
  determine the resulting policy before decrypting any constituent. The
  composition algebra operates over policy references in
  \texttt{FactOperational}, not over decrypted content. Otherwise the
  system decrypts Facts the requester is not allowed to see.

\item \textbf{Materialized vs.\ on-the-fly composition.} If a derived
  object is stable and shared (e.g., a published Narrative), it should
  carry a materialized composed policy. If it is reconstructed on
  demand, composition can be computed at access time.

\end{enumerate}

\textcolor{blue}{Need result type?}

\begin{verbatim}
pub enum DisclosureResult {
    Full(ResolvedEpisode),
    Partial {
        episode: ResolvedEpisode,   // built from authorized facts only
        excluded_count: usize,       // how many facts were excluded
        excluded_types: Vec<FactType>, // what categories, without identifying which
    },
    Denied,                          // no facts authorized for this purpose
}
\end{verbatim}

\section{Consent Revocation}
\label{sec:revocation}

A patient may revoke consent. Unless they are closing their account
entirely and requesting full data deletion, revocation must be granular:
scoped to specific Fact types, specific Episodes, specific borrowers, or
specific purposes.

\paragraph{Scope of revocation.}
Revocation prevents future access. It does not retroactively un-disclose
data already shared under a legitimate prior authorization. If a
borrower received decrypted Facts last week and the patient revokes
today, those Facts were legitimately accessed. The contract between the
patient and Phoros governs what the borrower must do with data already
in their possession (destroy, cease use, etc.), but the cryptographic
architecture cannot un-ring that bell.

\paragraph{Mechanism.}
Revocation invalidates the key material or access grants associated with
the revoked consent scope. If keys are scoped per-policy and the patient
revokes a policy, the system stops issuing keys under that policy.
Cached plaintexts in any operational system should have short TTLs and
be invalidated on revocation events.

\paragraph{Contractual framing.}
The initial contract between the patient and Phoros specifies timelines
and procedures for satisfying consent revocation requests. This includes
the window within which revocation takes effect, the obligations of
downstream borrowers, and the process for confirming deletion or
cessation of use.

\section{Provenance and Authorship}
\label{sec:provenance}

When a Fact records ``authored by Dr.~X,'' what binds that claim? In the
current design, the system verifies identity at the UI layer (e.g.,
authenticated session) and records authorship as a field within the
Fact. Once written, authorship is encrypted with all other fields of the
Fact and shares their protection.

This provides authentication but not non-repudiation. Dr.~X cannot
later prove they did not author the Fact, and Phoros cannot prove to a
third party that Dr.~X did, because the binding between identity
verification and the encrypted record is mediated entirely by Phoros's
systems.

If third-party verifiability becomes a requirement---for legal
proceedings, regulatory audits, or cross-institutional trust---the
architecture would need digital signatures on Facts at authoring time.
This is deferred.

\section{Summary}

The security architecture is organized around four commitments:

\begin{enumerate}

\item \textbf{Semantic alignment.} Protection boundaries follow FEN
  semantic units, not storage artifacts. Facts are the cryptographic
  unit. Episodes and Narratives inherit constraints from their
  constituents.

\item \textbf{Eager materialization.} FEN objects are created at
  ingestion. The object being encrypted and the object being governed
  are the same object.

\item \textbf{Transient plaintext.} Readable objects exist only in
  memory during authorized access. The system stores encrypted
  fragments and reconstructs semantic objects on demand through a
  policy-governed procedure.

\item \textbf{Operational separation.} \texttt{FactOperational}
  supports platform operations (indexing, discovery, policy evaluation)
  without decrypting payloads. Access to operational metadata is itself
  governed and audited.

\end{enumerate}

Consent composition follows intersection by default. Revocation is
granular and forward-looking. Key management and the signing mechanism
for non-repudiation are deferred to subsequent documents.

\end{document}